\section{Predictive Memory Localization}

PML connects a localized residual representation to a measured strength-indexed
intervention path and a held-out decision (Figure~\ref{fig:pml-pipeline}). Here,
``memory localization'' denotes a record-conditioned internal signal associated
with an answer relation, not a claim that the relation resides in one unique
physical component.

\subsection{Records, Probes, and Interventions}

Record $i$ contains target prompts $\mathcal{T}_i$, semantic-neighbor prompts
$\mathcal{N}_i$, and general-capability prompts $\mathcal{C}_i$. The first set
expresses the behavior to suppress or enhance; the latter two test whether
nearby knowledge or unrelated abilities are preserved. For prompt $x$ with
desired answer $y^+$ and contrast $y^-$, the answer margin is
\begin{equation}
    m(x)=\log p(y^+\mid x)-\log p(y^-\mid x).
\end{equation}
All outcomes are changes from the unperturbed margin, making paths comparable
across records with different baseline confidence.

At layer $\ell$, method $s$ constructs a unit-norm record-specific direction
$\mathbf{v}_{i\ell s}$. Intervention strength $\alpha$ modifies the residual
activation as
\begin{equation}
    \mathbf{h}'_{\ell}=\mathbf{h}_{\ell}+\alpha\mathbf{v}_{i\ell s}.
\end{equation}
Negative and positive coefficients test suppression and enhancement with the
same direction. We write $\Delta^T_i(\alpha)$ for signed target-margin movement
and $\Delta^N_i(\alpha),\Delta^C_i(\alpha)$ for neighbor and capability
movement; a sufficiently negative collateral change is damage.

\subsection{Random-Calibrated Intervention Paths}

Target, neighbor, and capability responses have different null scales, so
random directions define separate 95th-percentile thresholds $\tau_T$,
$\tau_N$, and $\tau_C$. A target effect crosses $\tau_T$ in the intended
direction, while neighbor or capability damage crosses the corresponding
negative threshold. A \emph{clean strength} produces a target effect while
crossing neither damage threshold at that same coefficient.

Across an ordered, finite strength set, these events form a measured-grid
intervention path. Target and collateral onset identify the first observed
crossing, and adjacent clean strengths form an observed clean region. We use
these grid-defined events as the path estimand.
The confirmatory prediction task focuses on four directly measured events at
held-out stronger coefficients:
\emph{Target-any}, semantic-neighbor damage,
capability damage, and \emph{Clean-any}. Percentages count
record--method--layer paths rather than prompts or individual strengths.

\subsection{Prediction Task}

PML forecasts later path outcomes from progressively stronger evidence.
Baseline belief $B$ and metadata $M$ are augmented with static localization
features $L$ and, for RFM, geometry features $G$. Low-dose response features
$R$ measure target and collateral movement at $\alpha=\pm0.1$. Because $R$
requires intervention, it is a cheap dynamic diagnostic rather than static
localization. The central comparisons test whether $L$ or $G$ improves on
$B+M$, and whether the strength-disjoint response $R$ forecasts outcomes at
$|\alpha|\in\{0.25,0.5\}$. A downstream policy then uses these forecasts to
select a coefficient or abstain.
